\section{Flag transitivity does not descend}
Theorem A in the introduction of Piedade and Tranchida,
\emph{Constructing new geometries: a generalized approach to
halving for hypertopes}, arXiv:2405.19050v1, asserts that the halved
geometry is flag-transitive if and only if the original geometry is.
We give a counterexample to the converse. The forward assertion in
the detailed Proposition 3.14 is unaffected. This example is not thin;
it is separate from the regular counterexample to residual connectivity.

\begin{lemma}\label{NT.twist}
Let $R$ be a simple graph on $V$, and let $g$ be an automorphism with
$g^2=1$. Suppose $x$ is not adjacent to $gx$ for every $x\in V$.
Define $Q$ by $x\sim_Q y$ if and only if $x\sim_R gy$.
Then $Q$ is simple and undirected. The canonical bipartite double
covers of $Q$ and $R$ are isomorphic by the identity on the first
shore and by $g$ on the second shore.
\end{lemma}
\begin{proof}
Symmetry follows from $x\sim_R gy$ if and only if $gx\sim_R y$,
and the hypothesis excludes a loop. In the double cover, $x$ on the
first shore is incident with $y$ on the second precisely when
$x\sim_R gy$. Thus the stated shore relabeling is the asserted
type-preserving isomorphism.
\end{proof}

Take $V$ to be the two-element subsets of $\{1,2,3,4,5\}$, let
$R$ have adjacency given by disjointness, and put $g=(12)$.
Every $X\in V$ meets $gX$: if $X$ contains exactly one of $1,2$,
its other element is fixed by $g$, and otherwise $X=gX$.
Lemma~\ref{NT.twist} therefore applies. It is useful to write
\[
 a=\{1,2\},\quad p_i=\{1,i\},\quad q_i=\{2,i\},\quad
 b_i=\{3,4,5\}\setminus\{i\}\qquad(i=3,4,5).
\]
The complete edge list of $Q$ is
\[
 \{p_i,p_j\},\ \{q_i,q_j\}\ (i<j),\qquad
 \{a,b_i\},\ \{b_i,p_i\},\ \{b_i,q_i\}\ (i=3,4,5).
\]
Indeed these are exactly the pairs $X,Y$ for which $X\cap gY$ is
empty, by separating the four displayed kinds of vertices. Thus
$Q$ is a connected cubic graph on ten vertices, with fifteen edges.
The $p_i$ and the $q_i$ each form a triangle; $a$ belongs to no
triangle, since its three neighbors are the pairwise nonadjacent
$b_i$. In particular $Q$ is nonbipartite and not vertex-transitive.

Define a rank-three incidence geometry $\Gamma$ as follows. Its
objects of type $0$ are the vertices of $Q$; its objects of type $1$
are the edges of $Q$, with the usual endpoint incidence. Its objects
of type $2$ are two distinct elements $z_+,z_-$, each incident with
every object of the other two types. Distinct objects of the same
type are not incident.

\begin{proposition}\label{NT.example}
The geometry $\Gamma$ is residually connected and satisfies the
source conditions B1 and B2. Its halving $\mathcal P(\Gamma)$ is
residually connected and flag-transitive, whereas $\Gamma$ is not
flag-transitive. It has $60$ chambers, as does its halving.
\end{proposition}
\begin{proof}
Every flag extends to a chamber $(v,e,z_\epsilon)$ with $v\in e$.
At $z_\epsilon$ the rank-two residue is the vertex-edge incidence
graph of $Q$, hence is connected. At a vertex it is $K_{3,2}$;
at an edge it is $K_{2,2}$. The whole incidence graph is connected
because either type-$2$ object is incident with all lower objects.
These are every residue of rank at least two, proving residual
connectivity. Simplicity of $Q$ is B1. Every type-$2$ object is
incident with both endpoints of every edge and with that edge, so
B2 holds. Flag transitivity would imply vertex transitivity of $Q$,
which its triangle distinction excludes. Finally there are
$15\cdot2\cdot2=60$ chambers.

We now identify the halving directly from its definition. Its types
$0$ and $1$ are two copies of $V$, with $X$ on the first shore
incident with $Y$ on the second when $X\cap gY=\varnothing$.
The shadow of either $z_\epsilon$ is the whole connected
nonbipartite graph $Q$. Its canonical double cover is connected,
so $z_\epsilon$ has one lifted object, incident with every object
of the first two types. Thus $\mathcal P(\Gamma)$ is the join of
the bipartite double cover of $Q$ with $\{z_+,z_-\}$.

For completeness, connectedness of that double cover follows without
an enumeration. In a connected nonbipartite graph, a walk between
two prescribed vertices can have either parity: traverse a path to
an odd cycle, go around that cycle, and return along the same path
to reverse the parity. Lifting these walks connects either choice
of sheet over either vertex. The residues in the new geometry are
therefore connected: at $z_\epsilon$ use this double cover, and at
either lower object the residue is $K_{3,2}$. These again exhaust
the rank-two residues.

After the second-shore relabeling $Y\mapsto gY$, the lower incidence
relation is simply disjointness of two-element subsets. The group
$S_5$ is transitive on ordered pairs of disjoint two-element sets:
map the first pair, the second pair, and their unique remaining
element separately. Consequently its action in the original labels
is
\[
 X_0\mapsto(\sigma X)_0,\qquad
 Y_1\mapsto(g\sigma gY)_1\qquad(\sigma\in S_5).
\]
It is transitive on the thirty incident lower pairs. Independently
interchanging $z_+$ and $z_-$ proves transitivity on all sixty
chambers. Since every flag extends to a chamber, this is flag
transitivity. The chamber count is $30\cdot2=60$.
\end{proof}

All nonempty higher shadows here equal $Q$ and are nonbipartite,
so the corrected parity condition is satisfied. The counterexample
therefore does not exploit the separate failure of unconditional
residual connectivity. Its only target is the reverse implication
in the introductory flag-transitivity equivalence. The converse
fails because an automorphism of the bipartite incidence geometry
may act differently on its two shores, and need not descend to an
automorphism of $Q$. The construction and proof use only elementary
graph and incidence arguments; no novelty or publication-priority
claim is made.
